% Part: computability
% Chapter: computability-theory
% Section: russells-paradox

\documentclass[../../../include/open-logic-section]{subfiles}

\begin{document}

\olfileid[pt]{cmp}{thy}{rus}
\olsection{Comparação com o paradoxo de Russell}

É instrutivo comparar e contrastar os argumentos desta seção com o
paradoxo de Russell:
\begin{enumerate}
\item Paradoxo de Russell: seja $S = \Setabs{x}{x \notin x}$. Então
  $S \in S$ se, e somente se, $S \notin S$, uma contradição.

  \emph{Conclusão:} Não existe tal conjunto~$S$. Supor a existência de
  um ``conjunto de todos os conjuntos'' é incompatível com os outros
  axiomas da teoria dos conjuntos.

\item Uma modificação do paradoxo de Russell: seja $F$ a ``função''
  do conjunto de todas as funções em $\{ 0, 1 \}$, definida por
  \[
  F(f) =
  \begin{cases}
    1 & \text{se $f$ pertence ao domínio de $f$, e $f(f) = 0$} \\
    0 & \text{caso contrário}
  \end{cases}
  \]
  Um argumento semelhante mostra que $F(F) = 0$ se, e somente se,
  $F(F) = 1$, uma contradição.

  \emph{Conclusão:} $F$ não é uma função. O ``conjunto de todas as
  funções'' é grande demais para ser o domínio de uma função.

\item O argumento de diagonalização: sejam $f_0$, $f_1$, \dots a
  enumeração das funções parcialmente computáveis, e seja
  $G \colon \Nat \to \{ 0, 1 \}$ definida por
  \[
  G(x) =
  \begin{cases}
    1 & \text{se $f_x(x)\downarrow = 0$} \\
    0 & \text{caso contrário}
  \end{cases}
  \]
  Se $G$ é computável, então ela é a função $f_k$ para algum
  $k$. Mas então $G(k) = 1$ se, e somente se, $G(k) = 0$, uma
  contradição.

  \emph{Conclusão:} $G$ não é computável. Observe que, de acordo com os
  axiomas da teoria dos conjuntos, $G$ ainda é uma função; não há
  paradoxo aqui, apenas um esclarecimento.
\end{enumerate}

Essa conversa sobre funções parciais, funções computáveis, funções
parcialmente computáveis e assim por diante pode ser confusa. O
conjunto de todas as funções parciais de $\Nat$ em $\Nat$ é uma grande
coleção de objetos. Algumas delas são totais, algumas são computáveis,
algumas são tanto totais quanto computáveis, e outras não são nem uma
coisa nem outra. Tenha em mente que, quando dizemos ``função'', por
padrão queremos dizer uma função total. Assim, temos:
\begin{enumerate}
\item funções computáveis
\item funções parcialmente computáveis que não são totais
\item funções que não são computáveis
\item funções parciais que não são nem totais nem computáveis
\end{enumerate}
Para organizar essas categorias, pode ser útil desenhar um quadrado
grande representando todas as funções parciais de $\Nat$ em $\Nat$ e
marcar, em seguida, duas regiões sobrepostas, correspondentes,
respectivamente, às funções totais e às funções parciais computáveis.
É um bom exercício verificar se você consegue descrever um objeto em
cada uma das regiões resultantes do diagrama.

\end{document}

